\documentclass[12pt,a4paper]{article}
\usepackage[utf8]{inputenc}
\usepackage[T1]{fontenc}
\usepackage[serbian]{babel}
\usepackage{amsmath,amssymb}
\usepackage[ruled,vlined,linesnumbered]{algorithm2e}
\usepackage{booktabs}
\usepackage{longtable}
\usepackage{geometry}
\usepackage{hyperref}
\usepackage{multirow}
\usepackage{array}
\usepackage{setspace}
\usepackage{parskip}
\usepackage{caption}

\geometry{left=2.5cm, right=2.5cm, top=2.5cm, bottom=2.5cm}
\setlength{\parindent}{0pt}
\setlength{\parskip}{6pt}

\hypersetup{colorlinks=true, linkcolor=black, urlcolor=blue, citecolor=black}

% ---------------------------------------------------------------
\title{\textbf{GRASP metaheuristika za dvostepeni transportni\\
problem sa fiksnim troškovima (TS-FCTP)}}
\author{Mateja Stojanović \\ Broj indeksa: mi251107}
\date{\today}
% ---------------------------------------------------------------

\begin{document}
\maketitle
\thispagestyle{empty}

% ---------------------------------------------------------------
\begin{abstract}
U ovom radu razmatra se Dvostepeni transportni problem sa fiksnim troškovima
(engl.\ \textit{Two-Stage Fixed-Charge Transportation Problem}, TS-FCTP),
koji modeluje distribuciju robe od dobavljača do distributivnih centara
i od distributivnih centara do kupaca uz varijabilne i fiksne troškove transporta.
Predložena je i implementirana metaheuristika zasnovana na GRASP pristupu
(engl.\ \textit{Greedy Randomized Adaptive Search Procedure}) sa nizom
proširenja: reaktivnim parametrom $\alpha$, lokalnom pretragom tipa VND
sa četiri susedstva, elitnim bazenom rešenja, prelaskom putanjom
(engl.\ \textit{path relinking}) i perturbacijom zasnovanom na
\textit{ruin-and-recreate} strategiji. Implementacija je dodatno
vektorizovana (NumPy) i paralelizovana (\texttt{multiprocessing}) na nivou
nezavisnih GRASP iteracija, čime je omogućeno izvršavanje algoritma više
puta ($k=3$) po instanci u razumnom vremenu i dobijanje statistički
smislenih pokazatelja ($\overline{gap}$, $\sigma$), umesto jednog
izvršavanja po instanci.

Algoritam je testiran na skupu od $60$ instanci malih, srednjih i velikih
dimenzija (po $20$ u svakoj grupi, $k=3$ izvršavanja po instanci) i
upoređen sa rešenjima dobijenim CPLEX solverom. Prosečno odstupanje od
CPLEX referentne vrednosti iznosi $\overline{gap}=9{,}24\%$
($\sigma=0{,}48\%$) na malim, $\overline{gap}=18{,}72\%$ ($\sigma=0{,}71\%$)
na srednjim i $\overline{gap}=29{,}19\%$ ($\sigma=0{,}92\%$) na velikim
instancama.
\end{abstract}
\vspace{50pt}
\textbf{Ključne reči:} dvostepeni transportni problem sa fiksnim troškovima,
GRASP, metaheuristika, lokalna pretraga, path relinking, vektorizacija,
paralelizacija, kombinatorna optimizacija.

\newpage
\tableofcontents
\newpage

% ===============================================================
\section{Opis problema i matematička formulacija}
% ===============================================================

\subsection{Definicija problema}

Dvostepeni transportni problem sa fiksnim troškovima (TS-FCTP) je problem
kombinatorne optimizacije koji se javlja u planiranju logističkih mreža.
Mreža se sastoji od tri sloja čvorova, povezanih u dve etape transporta:
\[
i \rightarrow j \rightarrow k,
\]
gde su:
\begin{itemize}
  \item $i \in I$ -- \textbf{dobavljači} (eng.\ \textit{suppliers}), sa raspoloživom ponudom $s_i$;
  \item $j \in J$ -- \textbf{distributivni centri} (DC);
  \item $k \in K$ -- \textbf{kupci} (eng.\ \textit{customers}), sa potražnjom $d_k$.
\end{itemize}

Za svaku vezu u mreži postoje varijabilni trošak (proporcionalan protoku)
i fiksni trošak koji se plaća jednom ukoliko se veza koristi, tj.\ ukoliko
je protok na njoj pozitivan. Cilj je minimizacija ukupnog troška transporta
uz poštovanje ograničenja kapaciteta dobavljača i ispunjenje potražnje kupaca.

\subsection{Skupovi i parametri}
\begin{itemize}
    \item $I$ -- skup dobavljača (izvora), $|I| = I$;
    \item $J$ -- skup distributivnih centara, $|J| = J$;
    \item $K$ -- skup kupaca, $|K| = K$;
    \item $s_i$ -- raspoloživa ponuda dobavljača $i$;
    \item $d_k$ -- potražnja kupca $k$;
    \item $b_{ij}$ -- varijabilni trošak transporta od dobavljača $i$ do centra $j$;
    \item $f_{ij}$ -- fiksni trošak aktiviranja veze $(i,j)$;
    \item $c_{jk}$ -- varijabilni trošak transporta od centra $j$ do kupca $k$;
    \item $g_{jk}$ -- fiksni trošak aktiviranja veze $(j,k)$.
\end{itemize}

\subsection{Promenljive}
\begin{itemize}
    \item $x_{ij} \ge 0$ -- količina koja se transportuje od dobavljača $i$ do centra $j$;
    \item $y_{jk} \ge 0$ -- količina koja se transportuje od centra $j$ do kupca $k$;
    \item $z_{ij} \in \{0,1\}$ -- indikator da li se koristi veza $(i,j)$;
    \item $w_{jk} \in \{0,1\}$ -- indikator da li se koristi veza $(j,k)$.
\end{itemize}

\subsection{Funkcija cilja}
Minimizuje se suma varijabilnih i fiksnih troškova u obe etape:
\[
\min \sum_{i \in I}\sum_{j \in J} \left( b_{ij}x_{ij} + f_{ij}z_{ij} \right)
+ \sum_{j \in J}\sum_{k \in K} \left( c_{jk}y_{jk} + g_{jk}w_{jk} \right).
\]

\subsection{Ograničenja}

\paragraph{Ograničenje ponude dobavljača}
\[
\sum_{j \in J} x_{ij} \le s_i, \quad \forall i \in I.
\]

\paragraph{Zadovoljenje potražnje kupaca}
\[
\sum_{j \in J} y_{jk} = d_k, \quad \forall k \in K.
\]

\paragraph{Balans protoka u distributivnim centrima}
Ukupni ulaz u centar mora biti jednak ukupnom izlazu -- DC ne skladišti
niti stvara robu, već je samo prosleđuje:
\[
\sum_{i \in I} x_{ij} - \sum_{k \in K} y_{jk} = 0, \quad \forall j \in J.
\]

\paragraph{Aktivacija fiksnih troškova (Big-M povezivanje)}
Ako se veza koristi (protok je pozitivan), odgovarajuća binarna promenljiva
mora biti $1$:
\[
x_{ij} \le s_i\, z_{ij}, \quad \forall i \in I, \forall j \in J,
\]
\[
y_{jk} \le d_k\, w_{jk}, \quad \forall j \in J, \forall k \in K.
\]

\subsection{Značenje funkcije cilja i uslova}
\begin{itemize}
    \item $b_{ij}x_{ij}$ i $c_{jk}y_{jk}$ su varijabilni troškovi proporcionalni
      protoku, dok su $f_{ij}z_{ij}$ i $g_{jk}w_{jk}$ fiksni troškovi koji se
      plaćaju jednom, ako se veza uopšte koristi.
    \item Ograničenja ponude obezbeđuju da nijedan dobavljač ne šalje više
      robe nego što ima na raspolaganju.
    \item Ograničenja potražnje garantuju da se zahtev svakog kupca tačno
      zadovolji.
    \item Balans u centrima znači da DC ne može praviti zalihe niti stvarati
      robu -- samo prosleđuje protok dalje.
    \item Big-M ograničenja povezuju kontinualne tokove sa binarnim
      indikatorima, čime se modeluje fiksni trošak aktivacije veze.
\end{itemize}

\subsection{Oblasti primene}

TS-FCTP se javlja u raznim oblastima:
\begin{itemize}
  \item \textbf{Upravljanje lancima snabdevanja} -- planiranje toka robe
    od fabrika do magacina i do maloprodajnih objekata;
  \item \textbf{Distribucija električne energije} -- alokacija energije od
    elektrana do transformatorskih stanica i do potrošača;
  \item \textbf{Telekomunikacije} -- rutiranje podataka kroz čvorove mreže;
  \item \textbf{Humanitarna logistika} -- distribucija pomoći od centralnih
    magacina do regionalnih centara i do krajnjih korisnika.
\end{itemize}

\subsection{Postojeći načini rešavanja}

Problem TS-FCTP pripada klasi NP-teških problema. Postojeći pristupi
uključuju:
\begin{itemize}
  \item \textbf{Tačni metodi}: MIP formulacija rešena komercijalnim
    solverima (npr.\ IBM ILOG CPLEX, Gurobi) uz primenu \textit{branch-and-bound}
    i \textit{branch-and-cut} algoritama. Za male instance postiže se
    optimalnost, ali za srednje i velike instance vreme rešavanja
    eksponencijalno raste.
  \item \textbf{Lagranževa relaksacija}: Relaksuju se ograničenja balansa
    ili kapaciteta uz Lagranževe multiplikatore koji se ažuriraju
    \textit{subgradijentnom metodom}. Dobijaju se dobre donje granice.
  \item \textbf{Tabu pretraga}: Iterativna lokalna pretraga sa listom
    zabrana za sprečavanje cikliranja. Efektivna za transportne probleme.
  \item \textbf{Genetski algoritmi / evolutivne metaheuristike}: Kombinuju
    više rešenja putem ukrštanja i mutacije.
  \item \textbf{GRASP}: Metaheuristika opisana u ovom radu, zasnovana na
    pohlepnoj nasumičnoj konstrukciji i lokalnoj pretrazi.
\end{itemize}

% ===============================================================
\section{Metaheuristika GRASP i njena adaptacija}
% ===============================================================

\subsection{Osnove GRASP-a}

GRASP (\textit{Greedy Randomized Adaptive Search Procedure}) je
iterativna metaheuristika koju su uveli Feo i Resende (1995).
Svaka iteracija se sastoji od dve faze:
\begin{enumerate}
  \item \textbf{Faza konstrukcije}: gradi se dopustivo rešenje pohlepno-nasumičnim
    pristupom. U svakom koraku formira se \textit{Restricted Candidate List}
    (RCL) -- skup kandidata čija je ocena u opsegu
    $[\text{min}, \text{min} + \alpha(\text{max}-\text{min})]$ --
    i nasumično se bira jedan element RCL-a.
    Parametar $\alpha \in [0,1]$ kontroliše stepen nasumičnosti:
    $\alpha=0$ odgovara čistoj pohlepnoj heuristici, a $\alpha=1$ nasumičnom izboru.
  \item \textbf{Faza lokalne pretrage}: dobijeno rešenje se poboljšava
    lokalnom pretragom u skupu susednih rešenja.
\end{enumerate}
Algoritam se ponavlja unapred zadati broj iteracija i pamti globalno
najbolje pronađeno rešenje.

\subsection{Implementacija}

Trenutna implementacija (folder \texttt{GRASP\_part\_parallel/} u repozitorijumu
projekta) sastoji se iz tri fajla: \texttt{utils.py} (učitavanje instance,
infrastruktura za multi-run eksperimente i izveštaj), \texttt{GRASP.py}
(ceo algoritam) i \texttt{main.py} (CLI ulaz). U nastavku je detaljno
opisana svaka klasa/funkcija iz \texttt{GRASP.py}, potom kompletan
pseudokod algoritma.

\subsubsection{Reprezentacija rešenja -- klasa \texttt{Solution}}

Rešenje se predstavlja parom matrica toka i skupom otvorenih DC-ova:
\begin{itemize}
  \item $x \in \mathbb{Z}_{\geq 0}^{I \times J}$ -- tokovi I. etape (dobavljač $\to$ DC);
  \item $y \in \mathbb{Z}_{\geq 0}^{J \times K}$ -- tokovi II. etape (DC $\to$ kupac);
  \item $\text{open\_dc} \subseteq \{1,\ldots,J\}$ -- skup otvorenih DC-ova
    (DC $j$ je otvoren ako $\sum_i x_{ij} > 0$ ili $\sum_k y_{jk} > 0$).
\end{itemize}

Klasa ima tri metode:
\begin{itemize}
  \item \texttt{copy()} -- duboka kopija rešenja (kopiraju se $x$, $y$, cena
    i \texttt{open\_dc}); koristi se jer se lokalna pretraga grana kroz
    probne (\textit{trial}) kopije koje se odbacuju ako ne donesu poboljšanje.
  \item \texttt{compute\_cost()} -- računa ukupnu cenu rešenja
    ($\sum b_{ij}x_{ij} + \sum c_{jk}y_{jk} + \sum f_{ij}[x_{ij}>0] + \sum g_{jk}[y_{jk}>0]$)
    i skup otvorenih DC-ova, u potpunosti vektorizovano preko NumPy-ja
    (\texttt{x.sum(axis=0) > 0 | y.sum(axis=1) > 0}, umesto petlje po $j$
    sa dva poziva sume po DC-u). Svaki poziv ove metode uvećava globalni
    brojač $eval$ (broj poziva funkcije cilja, videti odeljak o MA izveštaju).
  \item \texttt{is\_valid()} -- proverava da li je zadovoljena potražnja
    svakog kupca, da li nijedan dobavljač nije prekoračio kapacitet i da li
    važi balans toka u svakom DC-u.
\end{itemize}

Ključna arhitekturalna odluka: tokom pretrage se direktno modifikuju
isključivo $y$-tokovi (II. etapa); $x$-tokovi (I. etapa) se pri svakoj
promeni $y$-toka iznova rekonstruišu pohlepnom procedurom \textsc{Repair}.
Ovaj pristup smanjuje dimenzionalnost prostora pretrage jer algoritam
nikad eksplicitno ne upravlja $x$-promenljivama.

\subsubsection{Procedura \textsc{Repair}}

Za svaki DC $j$ sa pozitivnim ukupnim prilivom od kupaca ($\text{demand}_j
= \sum_k y_{jk}$), dobavljači se rangiraju po oceni koja kombinuje
varijabilni i amortizovani fiksni trošak,
$\text{score}_i = b_{ij} + f_{ij}/\text{demand}_j$, a zatim se tokovi
pohlepno dodeljuju u tom redosledu dok se zahtev DC-a ne zadovolji.
Ako neki DC ne može biti u potpunosti snabdeven (nedostatak ukupnog
kapaciteta), procedura vraća neuspeh.

\begin{algorithm}[H]
\caption{Repair$(sol, D)$}
\SetAlgoLined
$sol.x \leftarrow \mathbf{0}$\;
$\text{rem}\_s \leftarrow D.s.copy()$\;
\For{$j = 1$ \KwTo $J$}{
  $\text{demand}_j \leftarrow \sum_k y_{jk}$\;
  \lIf{$\text{demand}_j = 0$}{\textbf{continue}}
  Sortiraj dobavljače po $\text{score}_i = b_{ij} + f_{ij}/\text{demand}_j$\;
  $\text{rem} \leftarrow \text{demand}_j$\;
  \For{svakog $i$ u sortiranom redosledu}{
    $q \leftarrow \min(\text{rem}\_s[i],\, \text{rem})$\;
    $x_{ij} \mathrel{+}= q$;\quad
    $\text{rem}\_s[i] \mathrel{-}= q$;\quad
    $\text{rem} \mathrel{-}= q$\;
    \lIf{$\text{rem} = 0$}{\textbf{break}}
  }
  \lIf{$\text{rem} > 0$}{\Return \textit{neuspešno}}
}
\Return \textit{uspešno}
\end{algorithm}

Napomena o performansama: \textsc{Repair} je jedina preostala nevektorizovana
(sekvencijalna) tačka u lokalnoj pretrazi -- svaki DC deli tačno izbalansiran,
ograničen kapacitet dobavljača (ukupna ponuda je uvek jednaka ukupnoj
potražnji), pa bi inkrementalno prepravljanje samo pogođenih DC-ova moglo
promeniti da li \textsc{Repair} uspeva za pojedine kandidate -- taj rizik po
korektnost nije preuzet.

\subsubsection{Faza konstrukcije -- klasa \texttt{Construction}}

Kupci se obrađuju u opadajućem redosledu po potražnji. Za svakog kupca $k$,
u svakom koraku se formira RCL od parova $(i,j)$ sa ocenom:
\[
  \text{score}(i,j) = b_{ij} + c_{jk}
    + \frac{f_{ij}}{need}\cdot[x_{ij}=0]
    + \frac{g_{jk}}{need}\cdot[y_{jk}=0],
\]
gde je $need$ preostala potražnja kupca $k$ u tekućem koraku.
RCL sadrži sve parove čija je ocena
$\leq \text{score}_{\min} + \alpha(\text{score}_{\max} - \text{score}_{\min})$,
a iz njega se nasumično bira par kome se dodeljuje tok
$q = \min(\text{rem\_s}_i, need)$.

Cela $I \times J$ matrica ocena za tekućeg kupca gradi se i maskira
vektorizovano preko NumPy broadcast-a (umesto dvostruke petlje
\texttt{for i: for j:} koja generiše listu tuple-ova i sortira je), što je
za velike instance ($I,J$ reda stotina) dominantan doprinos ubrzanju
konstrukcije -- u merenjima na srednjoj instanci construction je posle
vektorizacije $>100\times$ brži po pozivu.

\begin{algorithm}[H]
\caption{GreedyRandomizedConstruction$(D, \alpha)$}
\SetAlgoLined
$sol \leftarrow$ prazno rešenje\;
\For{$k$ u opadajućem redosledu po $d_k$}{
  $need \leftarrow d_k$\;
  \While{$need > 0$}{
    Vektorizovano izračunaj $I \times J$ matricu ocena $\text{score}(i,j)$\;
    Maskiraj dobavljače sa $\text{rem}\_s[i] \le 0$ (ocena $= \infty$)\;
    $\text{score}_{\min} \leftarrow \min$, $\text{score}_{\max} \leftarrow \max$ (nad konačnim vrednostima)\;
    $\text{RCL} \leftarrow \{(i,j): \text{score}(i,j) \leq \text{score}_{\min} + \alpha(\text{score}_{\max}-\text{score}_{\min})\}$\;
    Nasumično izaberi $(i^*, j^*)$ iz RCL\;
    $q \leftarrow \min(\text{rem}\_s[i^*],\, need)$\;
    $x_{i^*j^*} \mathrel{+}= q$;\quad
    $y_{j^*k} \mathrel{+}= q$;\quad
    $\text{rem}\_s[i^*] \mathrel{-}= q$;\quad
    $need \mathrel{-}= q$\;
  }
}
Izračunaj cenu $sol$\;
\Return $sol$
\end{algorithm}

\subsubsection{Lokalna pretraga -- klasa \texttt{LocalSearch} (VND)}

Primenjuje se \textit{Variable Neighborhood Descent} (VND) sa četiri
susedstva, metodom \textit{first improvement} -- prvi pronađeni pomak koji
smanjuje cenu se prihvata i pretraga se vraća na prvo susedstvo.
Sve četiri metode koriste vektorizovanu (NumPy) evaluaciju kandidata:
umesto petlje \texttt{for k in range(K)} koja za svakog kupca
traži ko ga trenutno opslužuje (\texttt{next(j: y[j,k]>0)}), traži se
jednim pozivom \texttt{np.argmax(y > 0, axis=0)} za sve kupce odjednom.

\begin{enumerate}
  \item \textbf{Preraspoređivanje kupca} (\texttt{\_reassign\_customer}):
    za svakog kupca $k$ koji je dodeljen DC-u $j_{\text{from}}$,
    nasumično se isprobava do $10$ alternativnih DC-ova $j_{\text{to}}$
    (\texttt{random.sample}, umesto punog \texttt{shuffle} liste svih $J$
    DC-ova). Ako bi premeštanje smanjilo ukupnu cenu (jeftinija procena
    $c_{jk}d_k + g_{jk}$), premešta se ceo tok kupca (uz \textsc{Repair}
    i proveru stvarne cene).

  \item \textbf{Dodavanje DC-a} (\texttt{\_add\_dc}):
    isprobava se do $15$ zatvorenih DC-ova kao kandidati za otvaranje.
    Za svakog kandidata $j_{\text{new}}$, vektorizovano (numpy fensi-indeksiranje)
    se određuje skup kupaca kojima bi bilo jeftinije preko $j_{\text{new}}$,
    i svi se odjednom presmeravaju.

  \item \textbf{Uklanjanje DC-a} (\texttt{\_remove\_dc}):
    isprobava se do $5$ otvorenih DC-ova za zatvaranje; za svakog pogođenog
    kupca vektorizovano se bira najjeftinija preostala alternativa
    ($\arg\min$ po koloni matrice troška oblika $|\text{alts}| \times |\text{pogođeni kupci}|$).

  \item \textbf{Zamena DC-a} (\texttt{\_swap\_dc}):
    isprobava se do $30$ nasumičnih parova (otvoren DC, zatvoren DC) i
    izvršava zamena (svi kupci otvorenog DC-a presmeravaju se na zatvoreni)
    ako je poboljšanje moguće.
\end{enumerate}

Posle svakog uspešnog pomaka susedstvo se resetuje na prvo (klasičan VND
ciklus); ako nijedno od četiri susedstva ne poboljša rešenje
\texttt{max\_no\_improve}$=6$ puta uzastopno, pretraga se zaustavlja.

\begin{algorithm}[H]
\caption{VND$(D, sol, \text{max\_no\_improve}=6)$}
\SetAlgoLined
Izračunaj cenu $sol$\;
$\text{no\_improve} \leftarrow 0$\;
\While{$\text{no\_improve} < \text{max\_no\_improve}$}{
  \uIf{ReassignCustomer$(sol)$}{$\text{no\_improve} \leftarrow 0$; \textbf{continue}}
  \uElseIf{AddDC$(sol)$}{$\text{no\_improve} \leftarrow 0$; \textbf{continue}}
  \uElseIf{RemoveDC$(sol)$}{$\text{no\_improve} \leftarrow 0$; \textbf{continue}}
  \uElseIf{SwapDC$(sol)$}{$\text{no\_improve} \leftarrow 0$; \textbf{continue}}
  \Else{$\text{no\_improve} \mathrel{+}= 1$}
}
\Return $sol$
\end{algorithm}

\subsubsection{Elitni bazen -- klasa \texttt{ElitePool}}

Čuva se skup od do $12$ najkvalitetnijih i međusobno raznolikih rešenja.
Rešenje se ne dodaje u bazen ako u njemu već postoji rešenje sa istim
skupom otvorenih DC-ova (\texttt{open\_dc}) i sličnom cenom (razlika
$< 50$) -- na taj način bazen ne postaje prepun skoro identičnih rešenja.
Metoda \texttt{best()} vraća rešenje najniže cene, a \texttt{random()}
nasumično rešenje iz bazena; oba se koriste kao ulaz za \textit{path relinking}.

\subsubsection{Prelazak putanjom -- klasa \texttt{PathRelinking}}

Svakih $8$ iteracija, uzimaju se dva rešenja iz elitnog bazena (najbolje
i nasumično) i evaluiraju se \textbf{oba} kandidata za novi skup otvorenih
DC-ova: \textbf{unija} i \textbf{presek} njihovih skupova otvorenih DC-ova.
Za svaki od ta dva kandidatska skupa, svaki kupac se dodeljuje DC-u iz tog
skupa sa najnižom procenjenom cenom ($c_{jk}d_k + g_{jk}$), zatim se
primenjuje \textsc{Repair} i potom VND; zadržava se bolji od dva rezultata.

\begin{algorithm}[H]
\caption{PathRelinking$(a, b, D)$}
\SetAlgoLined
$best \leftarrow \text{null}$;\quad $best\_cost \leftarrow +\infty$\;
\For{$\text{dc\_set} \in \{a.\text{open\_dc} \cup b.\text{open\_dc},\; a.\text{open\_dc} \cap b.\text{open\_dc}\}$}{
  \lIf{$\text{dc\_set} = \varnothing$}{\textbf{continue}}
  $trial \leftarrow$ prazno rešenje\;
  \For{$k = 1$ \KwTo $K$}{
    $j^* \leftarrow \arg\min_{j \in \text{dc\_set}} \left(c_{jk}d_k + g_{jk}\right)$\;
    $y_{j^*k} \leftarrow d_k$\;
  }
  \If{\textsc{Repair}$(trial, D)$}{
    $trial \leftarrow$ VND$(D, trial)$\;
    \lIf{$trial.\text{cost} < best\_cost$}{$best \leftarrow trial$; $best\_cost \leftarrow trial.\text{cost}$}
  }
}
\Return $best$ (ili $a$, ako oba propadnu)
\end{algorithm}

\subsubsection{Perturbacija -- \textit{ruin-and-recreate} (shake)}

Ako se $20$ uzastopnih iteracija ne pronađe poboljšanje najboljeg rešenja,
primenjuje se perturbacija: DC-ovi se rangiraju po prosečnom varijabilnom
trošku po jedinici protoka, uklanja se $50\%$ najskupljih, a kupci koji su
im bili dodeljeni presmeravaju se na preostale (otvorene) DC-ove po
najnižoj proceni troška. Posle \textsc{Repair}-a i provere dopustivosti
primenjuje se VND; ako je novo rešenje bolje i dopustivo, prihvata se kao
novo najbolje. Brojač zastoja se resetuje posle svakog pokušaja shake-a
(bez obzira na uspeh), čime se sprečavaju uzastopni shake-ovi iz istog stanja.

\subsubsection{Reaktivni parametar $\alpha$}

Ako je prošlo više od $15$ iteracija bez poboljšanja najboljeg rešenja,
$\alpha$ se uzorkuje iz $[0{,}4,\, 0{,}7]$ radi diversifikacije pretrage
(veći RCL, više nasumičnosti). Inače, $\alpha \in [0{,}05,\, 0{,}5]$ radi
intenzifikacije (manji RCL, bliže pohlepnom izboru).

\subsubsection{Vektorizacija}

Profajliranjem (\texttt{cProfile}) je utvrđeno da su, pre optimizacije,
dva mesta dominirala vremenom izvršavanja: (1) konstrukcija -- $I \times J$
kandidata generisanih Python petljom po jedinici potražnje, i (2)
\texttt{compute\_cost()} -- računanje skupa otvorenih DC-ova Python petljom
od $J$ elemenata sa po dva poziva \texttt{np.sum()}. Oba mesta su zamenjena
NumPy vektorizovanim ekvivalentima (broadcast, boolean maske), 
bez promene rezultata pretrage -- ista RCL/$\alpha$ logika
i ista \textit{first-improvement} pravila, samo izračunata bržim putem.
Na srednjoj test instanci ovo je local search ubrzalo $\approx 1{,}75\times$
(sa 12{,}8~s na 7{,}3~s za 10 poziva construction+local-search).

\subsubsection{Paralelizacija -- batch nezavisnih iteracija}

GRASP petlja je inherentno sekvencijalna (elitni bazen, praćenje
najboljeg rešenja, \textit{path relinking} i shake zavise od deljenog
stanja), ali faza konstrukcije i lokalne pretrage jedne iteracije zavisi
samo od parametra $\alpha$ i nepromenjenih podataka instance -- potpuno je
nezavisna od ostalih iteracija. Zato je moguće paralelizovati batch od
$n_{\text{workers}}$ iteracija odjednom preko \texttt{multiprocessing.Pool}:

\begin{itemize}
  \item Podaci instance se svakom \textit{worker} procesu šalju tačno
    jednom, preko \texttt{initializer} funkcije \texttt{\_init\_worker}
    (ne pri svakom pozivu zadatka).
  \item Funkcija \texttt{\_construct\_and\_improve(alpha)} u worker procesu
    radi construction $+$ local search i vraća samo rezultat
    ($x$, $y$, cena, \texttt{open\_dc}, broj poziva funkcije cilja
    ) -- ne šalje nazad cost-matrice instance, koje su nepromenjene
    i već dostupne u worker procesu.
  \item Posle svakog batch-a, elitni bazen, praćenje najboljeg rešenja,
    \textit{path relinking} i shake primenjuju se \textbf{sekvencijalno},
    istim redosledom kao u serijskoj verziji -- jedina razlika je da su
    rešenja u batch-u već gotova (izračunata paralelno) pre te sekvencijalne
    obrade.
  \item Ako je $n_{\text{workers}} \le 1$, koristi se u potpunosti originalni
    serijski kod (\texttt{\_run\_serial}) -- ponašanje je tada identično
    algoritmu bez paralelizacije.
\end{itemize}

\begin{algorithm}[H]
\caption{GRASP\_Parallel$(D,\, T,\, n_{\text{workers}})$ -- razlika u odnosu na serijsku verziju}
\label{alg:parallel}
\SetAlgoLined
Pokreni \texttt{Pool} od $n_{\text{workers}}$ procesa, pošalji im $D$ jednom\;
$it \leftarrow 0$\;
\While{$it < T$}{
  $\text{batch} \leftarrow \min(n_{\text{workers}},\, T - it)$\;
  Generiši $\text{batch}$ vrednosti $\alpha$ (reaktivno, kao u serijskoj verziji)\;
  $\text{results} \leftarrow$ paralelno izračunaj \texttt{GreedyRandomizedConstruction}$+$VND za svaki $\alpha$ iz batch-a\;
  \For{svaki $sol$ iz $\text{results}$ (sekvencijalno, istim redosledom kao serijska GRASP petlja)}{
    Ažuriraj elitni bazen, najbolje rešenje, (svakih 8) path relinking, (posle 20 zastoja) shake\;
    $it \mathrel{+}= 1$\;
  }
}
Zatvori \texttt{Pool}\;
\Return najbolje pronađeno rešenje
\end{algorithm}

Broj \textit{worker} procesa se podrazumevano bira kao $\text{cpu\_count}()-1$
(ostavlja se jedno jezgro slobodno), a može se eksplicitno podesiti preko
CLI opcije \texttt{--workers}.

\subsubsection{Praćenje $t_i$ i $eval_i$ za MA izveštaj}

Da bi izveštaj (odeljak~3) bio potpun po specifikaciji, implementacija
prati, za svako izvršavanje:
\begin{itemize}
  \item $t_i$ -- \texttt{GRASP.run()} interno beleži \texttt{time.time()}
    u tačnom trenutku kada se najbolje rešenje poboljša (u glavnoj petlji,
    u \textit{path relinking}-u i u shake-u), i vraća tu vrednost kao deo
    rezultata. (Ranija verzija koda je ovde imala grešku -- $t_i$ je bio
    postavljen na ukupno vreme izvršavanja, bez stvarnog praćenja trenutka
    poboljšanja; to je ispravljeno.)
  \item $eval_i$ -- globalni brojač uvećan pri svakom pozivu
    \texttt{Solution.compute\_cost()}. U serijskom režimu čita se direktno
    na kraju izvršavanja; u paralelnom režimu svaki \textit{worker} proces
    ima sopstveni (nezavisan) brojač koji se resetuje na početku svakog
    zadatka i vraća kao deo rezultata, a glavni proces ga sabira sa
    sopstvenim pozivima (\textit{path relinking}, shake).
  \item $iter_i$ -- broj iteracija je fiksiran unapred (nema kriterijuma
    ranog zaustavljanja), pa je jednak zadatom broju iteracija $T$ za datu
    grupu instanci (videti odeljak~3.3).
  \item $cache_i$ -- u implementaciji ne postoji mehanizam kešfiguriranja
    (\textit{memoizacije}) poziva funkcije cilja, pa ovo (opciono) polje
    nije primenjivo i izostavljeno je iz izveštaja.
\end{itemize}

\subsubsection{Kompletan pseudokod GRASP algoritma}

\begin{algorithm}[H]
\caption{GRASP$(D,\, T)$}
\SetAlgoLined
\KwIn{Podaci instance $D$, broj iteracija $T$}
\KwOut{Najbolje pronađeno rešenje $sol^*$, vreme $t_i$ nalaska $sol^*$, broj evaluacija $eval$}
$sol^* \leftarrow \varnothing$;\quad
$\text{best\_cost} \leftarrow +\infty$;\quad
$\text{last\_imp} \leftarrow 0$;\quad
$t_i \leftarrow 0$;\quad
$E \leftarrow \varnothing$ (elitni bazen, veličina 12)\;
\For{$it = 0$ \KwTo $T-1$}{
  \eIf{$it - \text{last\_imp} > 15$}{
    $\alpha \leftarrow \mathcal{U}(0{,}4,\, 0{,}7)$\;
  }{
    $\alpha \leftarrow \mathcal{U}(0{,}05,\, 0{,}5)$\;
  }
  $sol \leftarrow$ GreedyRandomizedConstruction$(D, \alpha)$\;
  $sol \leftarrow$ VND$(D, sol)$\;
  $E.$add$(sol)$\;
  \If{$sol.\text{cost} < \text{best\_cost}$}{
    $sol^* \leftarrow sol$;\quad
    $\text{best\_cost} \leftarrow sol.\text{cost}$;\quad
    $\text{last\_imp} \leftarrow it$;\quad
    $t_i \leftarrow \text{time}() - \text{start}$\;
  }
  \If{$it \bmod 8 = 0$ \textbf{ i } $|E| \geq 2$}{
    $cand \leftarrow$ PathRelinking$(E.\text{best}(),\, E.\text{random}(),\, D)$\;
    \If{$cand.\text{cost} < \text{best\_cost}$}{
      $sol^* \leftarrow cand$;\quad
      $\text{best\_cost} \leftarrow cand.\text{cost}$;\quad
      $\text{last\_imp} \leftarrow it$;\quad
      $t_i \leftarrow \text{time}() - \text{start}$\;
    }
  }
  \If{$it - \text{last\_imp} > 20$}{
    $sol' \leftarrow$ RuinAndRecreate$(sol^*, D)$\;
    \If{\textsc{Repair}$(sol', D)$ i $sol'.\text{is\_valid}()$}{
      $sol' \leftarrow$ VND$(D, sol')$\;
      \If{$sol'.\text{is\_valid}()$ i $sol'.\text{cost} < \text{best\_cost}$}{
        $sol^* \leftarrow sol'$;\quad
        $\text{best\_cost} \leftarrow sol'.\text{cost}$;\quad
        $t_i \leftarrow \text{time}() - \text{start}$\;
      }
    }
    $\text{last\_imp} \leftarrow it$ \tcp*{sprečava uzastopne shake-ove}
  }
}
\Return $sol^*$, $t_i$, $eval$
\end{algorithm}

Paralelna verzija (Algoritam~\ref{alg:parallel}) izvršava telo petlje po $it$ na identičan
način, s tim što se \texttt{GreedyRandomizedConstruction}$+$VND za grupu od
$n_{\text{workers}}$ uzastopnih vrednosti $it$ računaju paralelno pre nego
što se sekvencijalni deo (ažuriranje $E$, $sol^*$, path relinking, shake)
primeni na svaki rezultat iz batch-a, tim redosledom.

% ===============================================================
\section{Rezultati testiranja i analiza}
% ===============================================================

\subsection{Test instance}

Za testiranje je korišćen skup od $60$ instanci podeljenih u tri grupe
(generisane koordinatnim modelom -- čvorovi nasumično u kvadratu
$[-400,400]^2$, varijabilni troškovi su Euklidske distance, fiksni troškovi
su višekratnici tih distanci, $r_f \in [10,15]$ za I. etapu i
$r_g \in [5,10]$ za II. etapu):
\begin{itemize}
  \item \textbf{Male} (small, 20 instanci): $I \in [4, 18]$, $J \in [8, 36]$, $K \in [16, 80]$.
  \item \textbf{Srednje} (medium, 20 instanci): $I \in [20, 60]$, $J \in [40, 120]$, $K \in [80, 275]$.
  \item \textbf{Velike} (large, 20 instanci): $I \in [70, 160]$, $J \in [140, 320]$, $K \in [280, 800]$.
\end{itemize}
Svaki dobavljač ima kapacitet $s_i = 80$; potražnja kupaca $d_k \in [10,30]$
je generisana nasumično, a zatim normalizovana tako da ukupna potražnja
tačno odgovara ukupnoj ponudi ($\sum_k d_k = \sum_i s_i$).

\subsection{Referentna rešenja (CPLEX)}

Referentna rešenja dobijena su CPLEX solverom (IBM ILOG CPLEX) korišćenjem
MIP formulacije iz odeljka~1, sa vremenskim limitom od $7200$ s i jednom niti.
Status \texttt{OPTI} označava da je CPLEX dokazao optimalnost; status
\texttt{FEAS} da je dostignut vremenski limit i vraćeno najbolje pronađeno
dopustivo rešenje (bez dokaza optimalnosti) -- to je i dalje referentna
vrednost korišćena za $Bestsol$, u skladu sa napomenom iz specifikacije
izveštaja ("optimalno ili najbolje poznato rešenje").


\subsection{Podešavanja GRASP-a, platforma i metodologija izveštaja}

\textit{Platforma za GRASP}: x86\_64, 16 CPU jezgara, Python 3.12.3,
NumPy 2.4.4. Algoritam je izvršavan paralelno sa
$n_{\text{workers}} = \text{cpu\_count}()-1 = 15$ worker procesa.

Svaka instanca je izvršena $k=3$ puta (seed $42 + 100r$ za $r$-to
izvršavanje, $r=0,1,2$), sa fiksnim brojem iteracija po grupi:
\begin{itemize}
  \item Male instance: $T = 500$;
  \item Srednje instance: $T = 200$;
  \item Velike instance: $T = 100$.
\end{itemize}

Za svako izvršavanje $i$ pamte se $sol_i$ (najbolja cena), $t_i$ (vreme do
prvog nalaska $sol_i$), $ttot_i$ (ukupno vreme izvršavanja) i $eval_i$ (broj
poziva funkcije cilja) -- videti odeljak~2.2.7. $Bestsol$ se određuje kao
CPLEX referentna vrednost ako je poznata i bolja od najboljeg MA rezultata,
u suprotnom kao najbolji MA rezultat kroz $k$ izvršavanja; zatim se za
svako $i$ računa
\[
  gap_i = 100 \cdot \frac{|sol_i - Bestsol|}{|Bestsol|}.
\]
U tabelama su prikazane srednje vrednosti $\bar t$, $\overline{ttot}$,
$\overline{eval}$, $\overline{gap}$ i standardna devijacija $\sigma$ (u \%),
računate preko $k=3$ izvršavanja po instanci.

\subsection{Male instance}

\begin{longtable}{r l r r r r r r r}
\caption{Rezultati za male instance ($T=500$, $k=3$)} \\
\toprule
\# & $I/J/K$ & Ref.\,(CPLEX) & Best MA & $\bar t$\,(s) & $\overline{ttot}$\,(s) & $\overline{eval}$\,($\times 10^3$) & $\overline{gap}$\,(\%) & $\sigma$\,(\%) \\
\midrule
\endfirsthead
\multicolumn{9}{c}{\tablename\ \thetable{} -- nastavak} \\
\toprule
\# & $I/J/K$ & Ref.\,(CPLEX) & Best MA & $\bar t$\,(s) & $\overline{ttot}$\,(s) & $\overline{eval}$\,($\times 10^3$) & $\overline{gap}$\,(\%) & $\sigma$\,(\%) \\
\midrule
\endhead
\midrule
\multicolumn{9}{r}{\textit{Nastavlja se na sledećoj strani}} \\
\endfoot
\bottomrule
\endlastfoot
1 & 4/8/16 & 150\,568 (OPTI) & 151\,577 & 0.9 & 3.9 & 135.6 & 0.67 & 0.00 \\
2 & 4/8/20 & 156\,022 (OPTI) & 164\,939 & 1.8 & 4.9 & 202.0 & 6.14 & 0.30 \\
3 & 4/10/20 & 161\,224 (OPTI) & 165\,831 & 4.0 & 8.8 & 305.5 & 2.87 & 0.01 \\
4 & 6/12/24 & 186\,537 (OPTI) & 195\,722 & 5.8 & 9.0 & 287.9 & 4.92 & 0.00 \\
5 & 6/12/30 & 213\,785 (OPTI) & 234\,361 & 2.1 & 14.0 & 493.0 & 10.56 & 0.82 \\
6 & 6/15/30 & 188\,954 (OPTI) & 206\,385 & 8.4 & 14.8 & 485.4 & 9.65 & 0.36 \\
7 & 8/16/32 & 231\,984 (OPTI) & 248\,473 & 10.0 & 13.3 & 401.7 & 7.17 & 0.05 \\
8 & 8/16/40 & 202\,467 (OPTI) & 219\,349 & 8.1 & 18.7 & 553.2 & 9.41 & 0.89 \\
9 & 8/20/40 & 210\,740 (OPTI) & 224\,763 & 20.0 & 22.7 & 612.3 & 7.09 & 0.36 \\
10 & 10/20/40 & 258\,236 (OPTI) & 275\,995 & 19.0 & 28.7 & 650.4 & 8.23 & 0.96 \\
11 & 10/20/50 & 286\,991 (OPTI) & 310\,636 & 17.7 & 31.2 & 735.3 & 8.67 & 0.38 \\
12 & 10/25/50 & 284\,643 (FEAS) & 314\,751 & 15.4 & 45.8 & 1179.3 & 11.08 & 0.62 \\
13 & 12/24/48 & 278\,044 (OPTI) & 316\,830 & 31.0 & 51.1 & 1011.4 & 14.26 & 0.26 \\
14 & 12/24/60 & 284\,648 (OPTI) & 328\,715 & 14.3 & 40.5 & 869.3 & 16.85 & 1.01 \\
15 & 12/30/60 & 315\,417 (OPTI) & 340\,537 & 34.3 & 65.6 & 1273.3 & 9.03 & 1.23 \\
16 & 14/28/56 & 349\,820 (FEAS) & 385\,816 & 37.8 & 65.7 & 1229.8 & 10.96 & 0.55 \\
17 & 14/28/70 & 368\,942 (FEAS) & 408\,612 & 65.5 & 82.1 & 1507.8 & 11.09 & 0.35 \\
18 & 16/32/64 & 349\,036 (FEAS) & 384\,333 & 39.4 & 72.6 & 1227.1 & 10.50 & 0.48 \\
19 & 16/32/80 & 351\,345 (FEAS) & 397\,712 & 40.2 & 88.3 & 1444.9 & 13.35 & 0.17 \\
20 & 18/36/72 & 400\,600 (FEAS) & 446\,026 & 25.4 & 93.9 & 1380.8 & 12.36 & 0.72 \\
\midrule
\multicolumn{4}{r}{\textbf{Prosek}} & 20.1 & 38.8 & 799.3 & \textbf{9.24} & \textbf{0.48} \\
\end{longtable}

\subsection{Srednje instance}

\begin{longtable}{r l r r r r r r r}
\caption{Rezultati za srednje instance ($T=200$, $k=3$)} \\
\toprule
\# & $I/J/K$ & Ref.\,(CPLEX) & Best MA & $\bar t$\,(s) & $\overline{ttot}$\,(s) & $\overline{eval}$\,($\times 10^3$) & $\overline{gap}$\,(\%) & $\sigma$\,(\%) \\
\midrule
\endfirsthead
\multicolumn{9}{c}{\tablename\ \thetable{} -- nastavak} \\
\toprule
\# & $I/J/K$ & Ref.\,(CPLEX) & Best MA & $\bar t$\,(s) & $\overline{ttot}$\,(s) & $\overline{eval}$\,($\times 10^3$) & $\overline{gap}$\,(\%) & $\sigma$\,(\%) \\
\midrule
\endhead
\midrule
\multicolumn{9}{r}{\textit{Nastavlja se na sledećoj strani}} \\
\endfoot
\bottomrule
\endlastfoot
1 & 20/40/80 & 383\,558 (FEAS) & 421\,275 & 26.8 & 41.5 & 558.1 & 10.95 & 1.01 \\
2 & 20/40/100 & 415\,983 (FEAS) & 470\,055 & 36.9 & 69.7 & 828.1 & 13.52 & 0.69 \\
3 & 20/50/100 & 400\,875 (FEAS) & 464\,645 & 48.9 & 78.3 & 839.9 & 17.32 & 1.03 \\
4 & 25/50/100 & 395\,660 (FEAS) & 464\,691 & 23.3 & 70.5 & 706.8 & 17.75 & 0.26 \\
5 & 25/50/125 & 474\,978 (FEAS) & 556\,570 & 63.4 & 106.8 & 1069.1 & 18.39 & 0.88 \\
6 & 25/60/120 & 447\,777 (FEAS) & 528\,696 & 107.2 & 160.4 & 1496.8 & 18.63 & 0.40 \\
7 & 30/60/120 & 474\,863 (FEAS) & 533\,621 & 58.3 & 96.5 & 820.8 & 13.84 & 1.10 \\
8 & 30/60/150 & 514\,195 (FEAS) & 619\,351 & 106.4 & 211.9 & 1663.8 & 21.09 & 0.77 \\
9 & 30/75/150 & 503\,494 (FEAS) & 580\,773 & 134.2 & 208.5 & 1413.1 & 16.21 & 0.62 \\
10 & 35/70/140 & 525\,261 (FEAS) & 619\,496 & 113.5 & 185.0 & 1416.3 & 18.24 & 0.22 \\
11 & 35/70/175 & 551\,427 (FEAS) & 647\,867 & 175.9 & 237.8 & 1622.1 & 18.53 & 0.81 \\
12 & 40/80/160 & 478\,570 (FEAS) & 544\,689 & 172.1 & 225.8 & 1352.2 & 14.94 & 0.82 \\
13 & 40/80/200 & 616\,913 (FEAS) & 764\,595 & 177.1 & 439.3 & 2832.9 & 24.16 & 0.15 \\
14 & 45/90/180 & 689\,812 (FEAS) & 829\,035 & 258.9 & 345.2 & 2039.4 & 21.20 & 0.84 \\
15 & 45/90/225 & 557\,928 (FEAS) & 657\,472 & 193.6 & 369.1 & 1956.9 & 20.05 & 1.73 \\
16 & 50/100/200 & 592\,227 (FEAS) & 721\,497 & 273.0 & 482.3 & 2396.1 & 22.40 & 0.41 \\
17 & 50/100/250 & 612\,725 (FEAS) & 733\,820 & 292.9 & 708.0 & 3534.2 & 21.36 & 1.22 \\
18 & 55/110/220 & 685\,116 (FEAS) & 858\,810 & 439.4 & 590.1 & 2723.6 & 25.57 & 0.16 \\
19 & 55/110/275 & 580\,226 (FEAS) & 686\,397 & 378.3 & 644.7 & 2766.0 & 19.17 & 0.65 \\
20 & 60/120/240 & 677\,683 (FEAS) & 816\,279 & 393.2 & 609.5 & 2433.1 & 21.11 & 0.49 \\
\midrule
\multicolumn{4}{r}{\textbf{Prosek}} & 173.7 & 294.0 & 1723.5 & \textbf{18.72} & \textbf{0.71} \\
\end{longtable}

\subsection{Velike instance}

\begin{longtable}{r l r r r r r r r}
\caption{Rezultati za velike instance ($T=100$, $k=3$)} \\
\toprule
\# & $I/J/K$ & Ref.\,(CPLEX) & Best MA & $\bar t$\,(s) & $\overline{ttot}$\,(s) & $\overline{eval}$\,($\times 10^3$) & $\overline{gap}$\,(\%) & $\sigma$\,(\%) \\
\midrule
\endfirsthead
\multicolumn{9}{c}{\tablename\ \thetable{} -- nastavak} \\
\toprule
\# & $I/J/K$ & Ref.\,(CPLEX) & Best MA & $\bar t$\,(s) & $\overline{ttot}$\,(s) & $\overline{eval}$\,($\times 10^3$) & $\overline{gap}$\,(\%) & $\sigma$\,(\%) \\
\midrule
\endhead
\midrule
\multicolumn{9}{r}{\textit{Nastavlja se na sledećoj strani}} \\
\endfoot
\bottomrule
\endlastfoot
1 & 70/140/280 & 741\,988 (FEAS) & 936\,363 & 296.9 & 508.4 & 1499.5 & 26.96 & 0.99 \\
2 & 70/140/350 & 825\,845 (FEAS) & 1\,036\,681 & 684.7 & 825.9 & 2600.4 & 27.41 & 1.35 \\
3 & 80/160/320 & 851\,981 (FEAS) & 1\,104\,103 & 452.1 & 844.6 & 2219.2 & 31.01 & 1.77 \\
4 & 80/160/400 & 807\,109 (FEAS) & 1\,013\,458 & 738.0 & 1117.9 & 2709.3 & 25.72 & 0.15 \\
5 & 90/180/360 & 801\,015 (FEAS) & 994\,840 & 289.8 & 1044.2 & 2261.4 & 24.44 & 0.25 \\
6 & 90/180/450 & 953\,624 (FEAS) & 1\,207\,206 & 1035.0 & 1456.3 & 3062.8 & 26.75 & 0.21 \\
7 & 100/200/400 & 1\,081\,573 (FEAS) & 1\,415\,328 & 690.1 & 1832.4 & 3231.7 & 31.88 & 0.90 \\
8 & 100/200/500 & 916\,397 (FEAS) & 1\,153\,410 & 1738.0 & 2921.7 & 4253.9 & 27.50 & 1.16 \\
9 & 110/220/440 & 871\,261 (FEAS) & 1\,136\,160 & 938.5 & 2107.5 & 3020.3 & 31.25 & 0.73 \\
10 & 110/220/550 & 945\,683 (FEAS) & 1\,189\,702 & 1466.8 & 2817.5 & 4391.2 & 26.49 & 0.63 \\
11 & 120/240/480 & 971\,579 (FEAS) & 1\,233\,209 & 1238.2 & 2292.6 & 3440.7 & 28.28 & 1.53 \\
12 & 120/240/600 & 1\,062\,116 (FEAS) & 1\,412\,545 & 2429.3 & 4425.9 & 6408.8 & 34.85 & 1.43 \\
13 & 130/260/520 & 958\,306 (FEAS) & 1\,207\,941 & 1935.1 & 4016.1 & 5183.2 & 27.06 & 1.23 \\
14 & 130/260/650 & 1\,017\,071 (FEAS) & 1\,308\,424 & 2691.6 & 4560.5 & 4931.7 & 29.80 & 1.01 \\
15 & 140/280/560 & 1\,015\,635 (FEAS) & 1\,274\,963 & 3843.3 & 4269.5 & 4807.6 & 26.03 & 0.57 \\
16 & 140/280/700 & 1\,020\,470 (FEAS) & 1\,260\,780 & 6278.5 & 7107.7 & 7383.9 & 24.58 & 1.32 \\
17 & 150/300/600 & 1\,418\,601 (FEAS) & 1\,955\,481 & 5316.6 & 7644.5 & 7542.3 & 39.03 & 0.86 \\
18 & 150/300/750 & 1\,125\,836 (FEAS) & 1\,518\,030 & 5697.4 & 9430.1 & 8442.3 & 36.51 & 1.19 \\
19 & 160/320/640 & 1\,191\,761 (FEAS) & 1\,519\,519 & 3602.3 & 6214.8 & 4998.4 & 28.05 & 0.41 \\
20 & 160/320/800 & 1\,178\,469 (FEAS) & 1\,529\,541 & 7921.3 & 10665.2 & 7277.5 & 30.29 & 0.67 \\
\midrule
\multicolumn{4}{r}{\textbf{Prosek}} & 2464.2 & 3805.2 & 4483.3 & \textbf{29.19} & \textbf{0.92} \\
\end{longtable}

\subsection{Analiza rezultata}

\subsubsection{Male instance}

Na malim instancama GRASP postiže prosečno odstupanje od $\overline{gap} =
9{,}24\%$ ($\sigma = 0{,}48\%$) u odnosu na CPLEX referentnu vrednost, uz
$k=3$ nezavisna izvršavanja po instanci -- standardna devijacija je mala
(uglavnom $< 1\%$), što ukazuje na konzistentno ponašanje algoritma za
fiksirani budžet od $T=500$ iteracija.
Najmanji gap postiže se na instanci~1 ($0{,}67\%$), najmanjih dimenzija
($4\times 8\times 16$), gde algoritam lako pronalazi skoro optimalno
rešenje. Najveći gap je na instanci~14 ($12 \times 24 \times 60$,
$16{,}85\%$), što ukazuje da porast broja kupaca po DC-u otežava lokalnu
pretragu da nadje globalni optimum u zadatom broju iteracija.
Prosečno vreme do najboljeg rešenja $\bar t$ raste sa dimenzijama instance
(od $0{,}9$~s na instanci~1 do $65{,}5$~s na instanci~17), dok je ukupno
vreme izvršavanja $\overline{ttot}$ svugde ispod $95$ sekundi zahvaljujući
paralelizaciji.

\subsubsection{Srednje instance}

Na srednjim instancama prosečno odstupanje iznosi $\overline{gap} =
18{,}72\%$ ($\sigma = 0{,}71\%$). Treba napomenuti da su sve CPLEX
referentne vrednosti za srednje instance označene statusom FEAS -- CPLEX
je dostigao vremenski limit od $7200$~s bez dokaza optimalnosti, pa
referentna vrednost nije nužno optimalna, već najbolje pronađeno dopustivo
rešenje kroz snažnu $B\&B$/\textit{cutting-plane} pretragu -- suštinski
drugačiji (i po vremenu mnogo skuplji) mehanizam od GRASP lokalne
pretrage sa $T=200$ iteracija.
Vreme izvršavanja raste sa dimenzijama, od $\overline{ttot}=41{,}5$~s
(instanca~1, $20\times40\times80$) do $\overline{ttot}=708{,}0$~s
(instanca~17, $50\times100\times250$), a broj evaluacija funkcije cilja
$\overline{eval}$ raste analogno, sa $558$ hiljada do preko $3{,}5$ miliona
poziva po izvršavanju.

\subsubsection{Velike instance}

Na velikim instancama prosečno odstupanje iznosi $\overline{gap} =
29{,}19\%$ ($\sigma = 0{,}92\%$), nastavljajući trend rasta uočen kod
malih i srednjih grupa ($9{,}24\%\to18{,}72\%\to29{,}19\%$) -- očekivano,
budući da je broj iteracija ($T=100$) najmanji od sve tri grupe, dok je
prostor pretrage ($I,J,K$ do $160,320,800$) najveći.
Najmanji gap je na instanci~5 ($90\times180\times360$, $24{,}44\%$), a
najveći na instanci~17 ($150\times300\times600$, $39{,}03\%$) --
zanimljivo, ne radi se nužno o najvećoj instanci u grupi (instanca~20,
$160\times320\times800$, ima gap $30{,}29\%$), što ukazuje da na ovim
dimenzijama i dalje postoji varijabilnost zavisna od konkretne strukture
instance, ne samo od njene veličine.
Vreme izvršavanja raste dramatično sa dimenzijama: $\overline{ttot}$ ide
od $508{,}4$~s (instanca~1) do preko $10\,665$~s (instanca~20, skoro
3 sata po izvršavanju), a broj evaluacija funkcije cilja $\overline{eval}$
raste do preko $8{,}4$ miliona poziva po izvršavanju (instanca~18) --
uprkos vektorizaciji i paralelizaciji, i dalje se radi o najzahtevnijoj
grupi instanci.

\subsubsection{Zajednički zaključci}

\begin{itemize}
  \item Odstupanje GRASP-a od CPLEX referentne vrednosti raste sa
    dimenzijom instance ($9{,}24\%\to18{,}72\%\to29{,}19\%$), što je
    očekivana posledica fiksnog broja iteracija po grupi ($500\to200\to100$)
    uz brzo rastuć prostor pretrage.
  \item Standardna devijacija $\sigma$ je kroz sve tri grupe relativno
    mala (prosečno $0{,}48$--$0{,}92\%$, pojedinačno najviše do
    $\approx 1{,}8\%$ na par instanci), što pokazuje da je algoritam, uz
    $k=3$ izvršavanja po instanci, dovoljno konzistentan -- razlike između
    pojedinačnih izvršavanja iste instance su znatno manje od razlika
    između instanci.
  \item Za srednje i velike instance (i deo malih) CPLEX referentne
    vrednosti same po sebi nisu dokazano optimalne (status FEAS), pa je
    stvarni gap GRASP-a prema (nepoznatom) optimumu verovatno nešto manji
    od prikazanog.
  \item Vektorizacija i paralelizacija (odeljci 2.2.9--2.2.10) su
    preduslov za to da su $k=3$ izvršavanja po svih $60$ instanci uopšte
    bila izvodljiva u razumnom vremenu -- pre tih izmena, jedno izvršavanje
    na srednjoj/velikoj instanci je trajalo od nekoliko minuta do više sati,
    a za velike instance ($T=100$, $k=3$, $20$ instanci) ukupno izvršavanje
    je i pored ubrzanja trajalo više desetina sati.
\end{itemize}

% ===============================================================
\section{Zaključak}
% ===============================================================

U ovom radu prezentovana je implementacija GRASP metaheuristike za
Dvostepeni transportni problem sa fiksnim troškovima (TS-FCTP). Algoritam
kombinuje pohlepno-nasumičnu konstrukciju sa reaktivnim parametrom $\alpha$,
VND lokalnu pretragu sa četiri susedstva, elitni bazen, path relinking
(unija i presek skupova otvorenih DC-ova) i shake perturbaciju.

Centralna arhitekturalna odluka -- direktna pretraga u prostoru $y$-tokova
uz rekonstrukciju $x$-tokova procedurom \textsc{Repair} -- efikasno smanjuje
dimenzionalnost problema. Nad tom osnovom, implementacija je dodatno
\textbf{vektorizovana} (NumPy, umesto Python petlji po $I,J,K$) i
\textbf{paralelizovana} (\texttt{multiprocessing}, batch nezavisnih
GRASP iteracija) -- što je omogućilo da se svaka instanca izvrši $k=3$
puta i dobiju statistički smislene ocene ($\overline{gap}$, $\sigma$),
umesto ranijeg pristupa sa $k=1$ zbog neizvodljivo dugih vremena
izvršavanja. Usput su ispravljene i dve greške u infrastrukturi za MA
izveštaj: $t_i$ (vreme do prvog nalaska najboljeg rešenja) se ranije nije
stvarno pratilo (bilo je jednako ukupnom vremenu), a referentna CPLEX
vrednost se zbog pogrešne putanje/poređenja imena fajla skoro nikad nije
učitavala, pa se $gap$ u praksi računao prema sopstvenom najboljem
rezultatu umesto prema CPLEX optimumu.

Testiranje na svih $60$ instanci ($k=3$ izvršavanja po instanci) pokazuje
da GRASP postiže dobra rešenja na malim instancama (prosečan gap
$9{,}24\%$, $\sigma=0{,}48\%$), sa umerenim padom kvaliteta na srednjim
instancama (gap $18{,}72\%$, $\sigma=0{,}71\%$) i daljim padom na velikim
instancama (gap $29{,}19\%$, $\sigma=0{,}92\%$), usled fiksnog i sve manjeg
broja iteracija ($500\to200\to100$) u odnosu na sve veći prostor pretrage
na krupnijim problemima. Kroz sve tri grupe standardna devijacija ostaje
relativno mala, što ukazuje na konzistentno ponašanje algoritma nezavisno
od dimenzije instance.

Mogući pravci daljeg unapređenja:
\begin{itemize}
  \item Povećati broj iteracija $T$ na srednjim i velikim instancama (npr.\
    dodatnim ubrzanjem \textsc{Repair}-a ili doniranjem više vremena
    izvršavanju) kako bi se smanjio uočeni gap na tim grupama;
  \item Uvesti mehanizam kešfiguriranja (\textit{caching}) funkcije cilja
    za ponovljena stanja, čime bi se omogućilo praćenje $cache_i$/$cache\%$
    iz specifikacije izveštaja i potencijalno smanjio broj stvarnih
    evaluacija;
  \item Inkrementalna (a ne potpuna) rekonstrukcija $x$-tokova u
    \textsc{Repair}-u za pogođene DC-ove, uz pažljivu proveru da ne naruši
    dopustivost usled tačno izbalansiranog kapaciteta;
  \item Ugradnja intenzivnijeg path relinking-a koji istražuje i
    međutačke na putu između dva elitna rešenja, ne samo krajnje uniju i
    presek;
\end{itemize}

\begin{thebibliography}{9}
\bibitem{tsfctp}
H.~I.~Calvete, C.~Gal\'e, J.~A.~Iranzo, P.~Toth,
\textit{A matheuristic for the two-stage fixed-charge transportation problem}.
\end{thebibliography}

\end{document}
